% GENERATED FILE. Edit canonical lessons, then run npm run generate:books.
% canonical-source-hash: fnv1a64:680a55b5ae9a7c34
% canonical-lessons: PT-C482-nuca, PT-C482-coxa, PT-C482-anca, PT-C482-calcanhar, PT-C482-umbigo, PT-R482-revisao-475-482

\chapter{Da nuca ao calcanhar}
\label{ch:pt-da-nuca-ao-calcanhar}
\hlchaptermodality{482}

\begin{chapteropening}
\textbf{By the end of this chapter:} \emph{I can name more parts of the body.}

Complete the last lesson of chapter 482: I can name more parts of the body.
\end{chapteropening}

\section[a nuca]{a nuca — nape (of the neck)}

\label{lesson:PT-C482-nuca}



\begin{quote}
Before the new one: say the Portuguese for nostril, then the Portuguese for gum.
\end{quote}

\subsection*{You'll want to know: a nuca}
\textbf{a nuca} — \textquotedblleft{}nape (of the neck)\textquotedblright{}.

\begin{quote}
O sol queimou-me a nuca.
\end{quote}

\begin{quote}
\emph{The sun burned the back of my neck.}
\end{quote}

The back of the neck; the front is \emph{a garganta}.

\begin{cousinweb}
\emph{nuca} comes from Medieval Latin \emph{nucha}, which is usually traced to an Arabic word for the spinal marrow.
\end{cousinweb}

\begin{grammarlens}[title={What you've built}]
One more word for telling, explaining and describing at B1.
\end{grammarlens}

\subsection*{Your turn}
\begin{itemize}
\raggedright
  \item \emph{Say it:} \emph{a nuca}
  \item \emph{Say it:} \emph{a nuca}, once more
  \item \emph{Say it:} a sentence of your own with \emph{a nuca}
  \item \emph{Recall:} say the Portuguese for nostril, then the Portuguese for gum, then say \emph{a nuca} again
\end{itemize}

\begin{tcolorbox}[breakable,colback=teal!4,colframe=teal!35!black,title={Before you move on}]
What does a nuca mean? (\textquotedblleft{}Nape (of the neck)\textquotedblright{}.) Say it once more.
\end{tcolorbox}
\section[a coxa]{a coxa — thigh}

\label{lesson:PT-C482-coxa}



\begin{quote}
Before the new one: say the Portuguese for gum, then the Portuguese for nape.
\end{quote}

\subsection*{You'll want to know: a coxa}
\textbf{a coxa} — \textquotedblleft{}thigh\textquotedblright{}.

\begin{quote}
Depois da corrida, doem-me as coxas.
\end{quote}

\begin{quote}
\emph{After the run, my thighs ache.}
\end{quote}

Of a chicken: \emph{a coxa de frango} — a chicken leg.

\begin{cousinweb}
\emph{coxa} is inherited, through Old Portuguese \emph{coixa} (EtymDB), from Latin \emph{coxa} \textquotedblleft{}hip\textquotedblright{} (Lewis \& Short); the meaning moved down to the thigh.
\end{cousinweb}

\begin{grammarlens}[title={What you've built}]
One more word for telling, explaining and describing at B1.
\end{grammarlens}

\subsection*{Your turn}
\begin{itemize}
\raggedright
  \item \emph{Say it:} \emph{a coxa}
  \item \emph{Say it:} \emph{a coxa}, once more
  \item \emph{Say it:} a sentence of your own with \emph{a coxa}
  \item \emph{Recall:} say the Portuguese for gum, then the Portuguese for nape, then say \emph{a coxa} again
\end{itemize}

\begin{tcolorbox}[breakable,colback=teal!4,colframe=teal!35!black,title={Before you move on}]
What does a coxa mean? (\textquotedblleft{}Thigh\textquotedblright{}.) Say it once more.
\end{tcolorbox}
\section[a anca]{a anca — hip}

\label{lesson:PT-C482-anca}



\begin{quote}
Before the new one: say the Portuguese for nape, then the Portuguese for thigh.
\end{quote}

\subsection*{You'll want to know: a anca}
\textbf{a anca} — \textquotedblleft{}hip\textquotedblright{}.

\begin{quote}
A avó partiu a anca numa queda.
\end{quote}

\begin{quote}
\emph{Grandma broke her hip in a fall.}
\end{quote}

In Brazil the usual word is \emph{o quadril}.

\begin{cousinweb}
\emph{anca} comes from Vulgar Latin \emph{hanca} (EtymDB), which was taken from a Germanic word for the hip.
\end{cousinweb}

\begin{grammarlens}[title={What you've built}]
One more word for telling, explaining and describing at B1.
\end{grammarlens}

\subsection*{Your turn}
\begin{itemize}
\raggedright
  \item \emph{Say it:} \emph{a anca}
  \item \emph{Say it:} \emph{a anca}, once more
  \item \emph{Say it:} a sentence of your own with \emph{a anca}
  \item \emph{Recall:} say the Portuguese for nape, then the Portuguese for thigh, then say \emph{a anca} again
\end{itemize}

\begin{tcolorbox}[breakable,colback=teal!4,colframe=teal!35!black,title={Before you move on}]
What does a anca mean? (\textquotedblleft{}Hip\textquotedblright{}.) Say it once more.
\end{tcolorbox}
\section[o calcanhar]{o calcanhar — heel}

\label{lesson:PT-C482-calcanhar}



\begin{quote}
Before the new one: say the Portuguese for thigh, then the Portuguese for hip.
\end{quote}

\subsection*{You'll want to know: o calcanhar}
\textbf{o calcanhar} — \textquotedblleft{}heel\textquotedblright{}.

\begin{quote}
Os sapatos novos fizeram-me uma ferida no calcanhar.
\end{quote}

\begin{quote}
\emph{The new shoes gave me a sore on my heel.}
\end{quote}

The heel of a shoe is \emph{o salto}.

\begin{cousinweb}
\emph{calcanhar} is usually traced to Vulgar Latin forms built on Latin \emph{calcāneum} \textquotedblleft{}heel\textquotedblright{} (Lewis \& Short), from \emph{calx} \textquotedblleft{}heel\textquotedblright{}.
\end{cousinweb}

\begin{grammarlens}[title={What you've built}]
One more word for telling, explaining and describing at B1.
\end{grammarlens}

\subsection*{Your turn}
\begin{itemize}
\raggedright
  \item \emph{Say it:} \emph{o calcanhar}
  \item \emph{Say it:} \emph{o calcanhar}, once more
  \item \emph{Say it:} a sentence of your own with \emph{o calcanhar}
  \item \emph{Recall:} say the Portuguese for thigh, then the Portuguese for hip, then say \emph{o calcanhar} again
\end{itemize}

\begin{tcolorbox}[breakable,colback=teal!4,colframe=teal!35!black,title={Before you move on}]
What does o calcanhar mean? (\textquotedblleft{}Heel\textquotedblright{}.) Say it once more.
\end{tcolorbox}
\section[o umbigo]{o umbigo — navel, belly button}

\label{lesson:PT-C482-umbigo}



\begin{quote}
Before the new one: say the Portuguese for hip, then the Portuguese for heel.
\end{quote}

\subsection*{You'll want to know: o umbigo}
\textbf{o umbigo} — \textquotedblleft{}navel, belly button\textquotedblright{}.

\begin{quote}
O bebé ri quando lhe tocam no umbigo.
\end{quote}

\begin{quote}
\emph{The baby laughs when someone touches his belly button.}
\end{quote}

\emph{Olhar para o próprio umbigo} — to think only of yourself.

\begin{cousinweb}
\emph{umbigo} is inherited, through Old Portuguese \emph{umbiigo}, from Latin \emph{umbilīcus} \textquotedblleft{}navel\textquotedblright{} (EtymDB).
\end{cousinweb}

\begin{grammarlens}[title={What you've built}]
One more word for telling, explaining and describing at B1.
\end{grammarlens}

\subsection*{Your turn}
\begin{itemize}
\raggedright
  \item \emph{Say it:} \emph{o umbigo}
  \item \emph{Say it:} \emph{o umbigo}, once more
  \item \emph{Say it:} a sentence of your own with \emph{o umbigo}
  \item \emph{Recall:} say the Portuguese for hip, then the Portuguese for heel, then say \emph{o umbigo} again
\end{itemize}

\begin{tcolorbox}[breakable,colback=teal!4,colframe=teal!35!black,title={Before you move on}]
What does o umbigo mean? (\textquotedblleft{}Navel, belly button\textquotedblright{}.) Say it once more.
\end{tcolorbox}
\section[(dialogue)]{Revisão dos capítulos 475–482}

\label{lesson:PT-R482-revisao-475-482}



\begin{quote}
Say the Portuguese for nervousness and to irritate.
\end{quote}

\subsection*{Your turn}
Say these aloud:

\begin{itemize}
\raggedright
  \item nervousness, to irritate, impatient, to calm, relief
  \item disappointment, to disappoint, discouraged, to console, to get something off your chest
  \item crying, to sob, to sigh, sigh, burst of laughter
  \item nostalgia, melancholy, loneliness, lonely, grief
  \item to hurt, hurt, grudge, to forgive, forgiveness
  \item to regret, regret, remorse, to embarrass, ashamed
  \item eyebrow, eyelash, eyelid, nostril, gum
  \item nape, thigh, hip, heel, navel
\end{itemize}

\begin{tcolorbox}[breakable,colback=teal!4,colframe=teal!35!black,title={Before you move on}]
Nervousness? (\textbf{o nervosismo}.) To irritate? (\textbf{irritar}.) Impatient? (\textbf{impaciente}.) To calm? (\textbf{acalmar}.) Relief? (\textbf{o alívio}.) Disappointment? (\textbf{a desilusão}.) To disappoint? (\textbf{desiludir}.) Discouraged? (\textbf{desanimado}.) To console? (\textbf{consolar}.) To get something off your chest? (\textbf{desabafar}.) Crying? (\textbf{o choro}.) To sob? (\textbf{soluçar}.) To sigh? (\textbf{suspirar}.) Sigh? (\textbf{o suspiro}.) Burst of laughter? (\textbf{a gargalhada}.) Nostalgia? (\textbf{a nostalgia}.) Melancholy? (\textbf{a melancolia}.) Loneliness? (\textbf{a solidão}.) Lonely? (\textbf{solitário}.) Grief? (\textbf{o desgosto}.) To hurt? (\textbf{magoar}.) Hurt? (\textbf{a mágoa}.) Grudge? (\textbf{o rancor}.) To forgive? (\textbf{perdoar}.) Forgiveness? (\textbf{o perdão}.) To regret? (\textbf{arrepender}.) Regret? (\textbf{o arrependimento}.) Remorse? (\textbf{o remorso}.) To embarrass? (\textbf{envergonhar}.) Ashamed? (\textbf{envergonhado}.) Eyebrow? (\textbf{a sobrancelha}.) Eyelash? (\textbf{a pestana}.) Eyelid? (\textbf{a pálpebra}.) Nostril? (\textbf{a narina}.) Gum? (\textbf{a gengiva}.) Nape? (\textbf{a nuca}.) Thigh? (\textbf{a coxa}.) Hip? (\textbf{a anca}.) Heel? (\textbf{o calcanhar}.) Navel? (\textbf{o umbigo}.)
\end{tcolorbox}
